\section*{Bab \ref{ch:b2:fluid-kinematics} --- Kinematika Fluida}

\begin{solution}{exo:b2:fluid-kinematics:1}
(a) $\ell/L \sim 7 \times 10^{-8}$: kontinum. (b) $0.07$: ambang --- hukum
kontinumnya mulai gagal (selip di dindingnya). (c) $\ell = 70 \times 10^{-9}
/3 \times 10^{-7} = \qty{0.23}{m}$, $\ell/L = 0.23$: bukan kontinum; gasnya
adalah hujan molekul yang saling bebas (rezim bebas molekuler).
\end{solution}

\begin{solution}{exo:b2:fluid-kinematics:2}
$a = v\,\dd v/\dd x = v_0^2(1 + x/L)/L$ (dari $v_0^2/L$ sampai $2v_0^2/L$).
$\dd x/\dd t = v_0(1 + x/L)$: $t = (L/v_0)\ln2 = 0.69\,L/v_0$, antara
$L/2v_0$ dan $L/v_0$.
\end{solution}

\begin{solution}{exo:b2:fluid-kinematics:3}
$D_V = \int_0^Rv_{\max}(1 - r^2/R^2)\,2\pi r\,\dd r = \tfrac12\pi R^2v_{\max}$;
$\langle v\rangle = v_{\max}/2 = \qty{0.20}{m/s}$; $D_V = \qty{1.6e-5}{m^3/s} =
\qty{16}{mL/s}$.
\end{solution}

\begin{solution}{exo:b2:fluid-kinematics:4}
(a) $\operatorname{div} = 0$, $\vect\omega = \vect 0$: tak termampatkan,
irotasional. (b) $\operatorname{div} = 0$, $\omega_z = 2\Omega$: rotasi
tegar. (c) $\operatorname{div} = 2a \ne 0$, irotasional (sebuah sumber
bidang). (d) $\operatorname{div} = 0$, $\omega_z = -k$: geser. (e)
$\operatorname{div} = b$: tak termampatkan jika dan hanya jika $b = 0$; irotasional untuk
setiap $b$.
\end{solution}

\begin{solution}{exo:b2:fluid-kinematics:5}
$D_V = \qty{0.20}{L/s}$. Selang $S = \qty{1.13}{cm^2}$: $v = \qty{1.8}{m/s}$; corong
$S = \qty{7.1}{mm^2}$: $v = \qty{28}{m/s}$. Percepatan rata-rata $\Delta(v^2/2)/L
= (28^2 - 1.8^2)/0.08 \approx \qty{1e4}{m/s^2} \approx 1000\,g$. Waktunya $\approx
\int\dd x/v$: beberapa milisekon (\qty{4}{ms} untuk profil kelajuan linear).
\end{solution}

\begin{solution}{exo:b2:fluid-kinematics:6}
(a) $\dd y/\dd x = (V/U)\cos\omega t$: garis lurus berkemiringan
$(V/U)\cos\omega t$, semuanya sejajar, memiring menurut waktu. (b) $x = Ut$, $y =
(V/\omega)\sin\omega t$: sinusoida $y = (V/\omega)\sin(\omega x/U)$. (c) Pada
$t = 0$ garis arusnya berada pada \qty{45}{\degree} dan \omterm{def:b2:fluid-kinematics:euler}{garis lintasannya} berangkat
dari titik asal pada \qty{45}{\degree}; pada $t = \pi/2\omega$ garis arusnya
mendatar sementara \omterm{def:b2:fluid-kinematics:euler}{garis lintasannya} sinusoida tetap. Keduanya berbeda sebab
alirannya tak tunak: \omterm{def:b2:fluid-kinematics:euler}{garis arus} adalah sebuah potret, \omterm{def:b2:fluid-kinematics:euler}{garis lintasan} sebuah riwayat.
(d) $\vect a = \partial_t\vect v = (0, -V\omega\sin\omega t)$ --- medannya
seragam, jadi suku konvektifnya lenyap.
\end{solution}

\begin{solution}{exo:b2:fluid-kinematics:7}
(a) $D_V = 50 \times 2 \times 1.2 = \qty{120}{m^3/s}$. (b) $120/60 = \qty{2.0}{m/s}$.
(c) $150/150 = \qty{1.0}{m/s}$. (d) $\rho \propto 1/T$ pada $P$ tetap menjadi separuh;
$\rho vS$ kekal: kelajuannya menjadi dua kali lipat.
\end{solution}

\begin{solution}{exo:b2:fluid-kinematics:8}
(a) $\omega_z = \frac1r\frac{\dd(rv_\theta)}{\dd r}$: $2\Omega$ di dalam, $0$
di luar. (b) Di dalam $\Gamma = 2\pi r\cdot\Omega r = 2\Omega\cdot\pi r^2$, yaitu fluks
vortisitasnya; di luar $\Gamma = 2\pi\Omega a^2$, tetap, yaitu fluks
seluruh intinya. (c) $\Omega = v_{\max}/a = \qty{1.4}{rad/s}$; $\Gamma = 2\pi av_{\max}
= \qty{2.2e4}{m^2/s}$; di \qty{500}{m}: $\Omega a^2/r = \qty{7}{m/s}$. (d)
Intinya berputar seperti \omterm{def:b2:rigid-body-mechanics:solid}{benda tegar}; di luarnya alirannya irotasional: partikelnya
mengelilingi sumbunya tanpa berputar terhadap dirinya sendiri. Vortisitasnya ---
rotasinya --- bersemayam di intinya.
\end{solution}

\begin{solution}{exo:b2:fluid-kinematics:9}
(a) $\varphi = \tfrac12k(x^2 - y^2)$, $\Delta\varphi = k - k = 0$. (b) $\dd x/
kx = \dd y/(-ky)$: $xy =$ tetap; $\partial_y\psi = kx = v_x$, $-\partial_x\psi =
-ky = v_y$. (c) $\vect a = (v_x\partial_x + v_y\partial_y)\vect v = k^2(x, y) =
\operatorname{\vect{grad}}\bigl(\tfrac12k^2(x^2 + y^2)\bigr) = \operatorname{\vect{grad}}
(v^2/2)$: dalam \omterm{def:b2:fluid-kinematics:vorticity}{aliran irotasional} tunak percepatannya adalah
gradien energi kinetik tiap satuan massa; ia menjauhi
titik stagnasinya --- sebuah partikel diperlambat ketika mendekati dindingnya
sepanjang $y$ dan dipercepat menjauh sepanjang $x$. (d) $\dot x = kx$, $\dot y = -ky$:
$x = x_0\eu^{kt}$, $y = y_0\eu^{-kt}$ ($xy$ tetap).
\end{solution}

\begin{solution}{exo:b2:fluid-kinematics:10}
(a) $\dd x/\dd t = x/(t + \tau)$ memberi $x = x_0(1 + t/\tau)$, $\ddot x = 0$;
\omterm{thm:b2:fluid-kinematics:material}{turunan partikelnya}: $\partial_tv + v\partial_xv = -x/(t + \tau)^2 + x/(t + \tau)^2
= 0$. (b) $\partial_t\rho + \partial_x(\rho v) = \dot\rho + \rho/(t + \tau) = 0$: $\rho =
\rho_0\tau/(t + \tau)$. (c) Jarak antara dua partikel tumbuh sebagai $(1 + t/\tau)$
dan massa di antaranya tetap: $\rho \propto 1/(1 + t/\tau)$. (d) $v =
Hx$ dengan $H = 1/(t + \tau)$: sebuah hukum Hubble, pemuaian bebas (tanpa
percepatan), massa jenisnya turun sebagai kebalikan faktor skalanya.
\end{solution}

\begin{solution}{exo:b2:fluid-kinematics:11}
(a) $\vect v = \operatorname{\vect{grad}}\varphi = (U + qx/2\pi r^2, qy/2\pi r^2)$;
sumbernya mempunyai $v_r = q/2\pi r$, $\operatorname{div} = \frac1r\dd(rv_r)/\dd r
= 0$ dan tak bergantung pada $\theta$, jadi $\omega_z = 0$; jumlah medan semacam itu
tak termampatkan dan irotasional pula. (b) Pada sumbu $x$ negatif,
$U - q/2\pi r = 0$: $x_s = -q/2\pi U$. (c) Di titik stagnasinya $\theta =
\pi$, $y = 0$: $\psi = q/2$; \omterm{def:b2:fluid-kinematics:euler}{garis arus} $Ur\sin\theta + q\theta/2\pi = q/2$,
yakni $r\sin\theta = (q/2\pi U)(\pi - \theta)$ --- kurva yang berangkat dari $x_s$ dan
membuka ke hilir. (d) $\theta \to 0$: $y \to q/2U$; $\theta \to 2\pi$ (cabang
bawah): $y \to -q/2U$. Fluida di dalam \omterm{def:b2:fluid-kinematics:euler}{garis arus} itu datang dari
sumbernya; di luarnya, arus seragamnya mengalir mengelilingi setengah benda
berlebar asimtotik $q/U$.
\end{solution}

\begin{solution}{exo:b2:fluid-kinematics:12}
(a) $\operatorname{div}\vect v = \frac1r\partial_r(-q/2\pi) = 0$; $\omega_z =
\frac1r\partial_r(\Gamma/2\pi) = 0$. (b) $\dd r/v_r = r\,\dd\theta/v_\theta$ memberi
$\dd r/r = -(q/\Gamma)\dd\theta$. (c) $a_r = v_r\partial_rv_r - v_\theta^2/r =
-(q^2 + \Gamma^2)/4\pi^2r^3 = -v^2/r$; $a_\theta = v_r\partial_rv_\theta + v_rv_\theta/r
= 0$. (d) $\dd r/\dd t = -q/2\pi r$: $t = \pi(r_0^2 - a^2)/q$; putarannya
$= (\Gamma/2\pi q)\ln(r_0/a)$. (e) $t = \qty{520}{s} \approx \qty{9}{min}$; $10$
putaran.
\end{solution}

\begin{solution}{pb:b2:fluid-kinematics:1}
\textbf{1.} Lewat tabung berjari-jari $r$: $2\pi rh|v_r| = D_V$, jadi
$v_r = -q/2\pi r$, $q = D_V/h = \qty{1.5e-3}{m^2/s}$.

\textbf{2.} $rv_\theta = 0.5 \times 0.01 = \qty{5e-3}{m^2/s}$: $v_\theta = \Gamma/2\pi r$
dengan $\Gamma = \qty{3.1e-2}{m^2/s}$.

\textbf{3.} $\operatorname{div}\vect v = \frac1r\partial_r(rv_r) = 0$, $\omega_z =
\frac1r\partial_r(rv_\theta) = 0$.

\textbf{4.} $\dd r/r = -(q/\Gamma)\dd\theta$: $r = r_0\eu^{-q\theta/\Gamma}$; putarannya
$= (\Gamma/2\pi q)\ln(r_0/a) = 3.3 \times 3.2 \approx 11$.

\textbf{5.} $t = \pi(r_0^2 - a^2)/q = \qty{520}{s} \approx \qty{9}{min}$; $v_\theta(a) =
\qty{0.25}{m/s}$.

\textbf{6.} $\vect a = -(v^2/r)\vect e_r$; di $r = a$: $v^2 = 0.25^2 + 0.012^2$,
$a = \qty{3.1}{m/s^2} \approx 0.3g$.

\textbf{7.} $\Omega_\oplus\sin\lambda\,r = 5 \times 10^{-5} \times 0.5 = \qty{2.5e-5}{m/s}$,
empat ratus kali lebih kecil daripada adukannya: belahan buminya tidak
menentukan apa pun dalam bak biasa (ia menentukan dalam tangki yang dibiarkan diam sehari).

\textbf{8.} $v_\theta = \Omega r$ ($r < a$), $\Omega a^2/r$ ($r > a$); $\Omega = 80/60
= \qty{1.33}{rad/s}$, $\Gamma = 2\pi av_{\max} = \qty{3.0e4}{m^2/s}$.

\textbf{9.} $\omega_z = 2\Omega = \qty{2.7}{s^{-1}}$ di intinya, $0$ di luarnya;
$v_\theta$ naik linear sampai \qty{80}{m/s} lalu turun sebagai $1/r$; $\omega_z$ berbentuk
tangga.

\textbf{10.} $a = v_{\max}^2/a = 6400/60 = \qty{107}{m/s^2} \approx 11g$, ke dalam.

\textbf{11.} Di luarnya vortisitasnya nol: unsur kecil bertranslasi
sepanjang lingkarannya tanpa berputar terhadap dirinya sendiri --- kincirnya
menjaga arah hadapnya. Di dalam intinya (rotasi tegar) ia ikut berputar bersama
fluidanya, sekali tiap satu kelilingan.

\textbf{12.} $E_{\text{inti}} = \int_0^a\tfrac12\rho\Omega^2r^2\,2\pi r\,\dd r =
\tfrac14\pi\rho\Omega^2a^4 = \qty{2.2e7}{J/m}$.

\textbf{13.} $E_{\text{luar}} = \int_a^R\tfrac12\rho(\Omega a^2/r)^2\,2\pi r\,\dd r =
\pi\rho\Omega^2a^4\ln(R/a) = 8.7 \times 10^7 \times 4.4 = \qty{3.8e8}{J/m}$;
integrannya $\rho v^2r \propto 1/r$, dari situlah logaritmanya.

\textbf{14.} $\qty{4.0e8}{J/m} \times \qty{1000}{m} = \qty{4e11}{J}$: sekitar
\qty{100}{t} TNT.

\textbf{15.} Aliran masuk $2\pi a \times 300 \times 5 = \qty{5.7e5}{m^3/s}$ lewat
$\pi a^2 = \qty{1.1e4}{m^2}$: $v_z = \qty{50}{m/s}$.

\textbf{16.} $\omega = 2v_{\max}/a = \qty{2.5e-3}{s^{-1}} = 50f$; $\Gamma = 2\pi av_{\max}
= \qty{1.3e7}{m^2/s}$.

\textbf{17.} $r_0\cdot fr_0/2 = r(v_\theta + fr/2)$: $v_\theta = f(r_0^2 - r^2)/2r$;
\qty{60}{m/s} di \qty{100}{km}, \qty{155}{m/s} di \qty{40}{km}.

\textbf{18.} $f > 0$ di belahan bumi utara memberi $v_\theta > 0$:
berlawanan jarum jam dilihat dari atas (siklonik). Atas \qty{500}{km} rotasi
planetnya adalah rotasi awal yang menguasai; dalam bak mandi
adukannyalah yang menguasai.

\textbf{19.} Gesekan dengan laut mengambil momentum sudutnya, dan
gradien tekanannya tak dapat mengimbangi percepatan sentripetal cincin
yang kian cepat: udaranya naik sebelum mencapai sumbunya, di dinding
matanya, sehingga menyisakan mata yang tenang.

\textbf{20.} $\rho\,2\pi rHv = 1.2 \times 2\pi \times 5 \times 10^5 \times 10^3 \times 5 =
\qty{1.9e10}{kg/s}$.

\textbf{21.} $\vect v(M, t) = -U\vect e_x + \vect v_0(M + Ut\,\vect e_x)$: polanya
terbawa arus, medannya tak tunak dalam kerangka tanah.
Translasinya menambah pusarannya di sisi tempat $v_\theta$ menunjuk
ke barat --- sisi utara (kanan) lintasan ke baratnya.

\textbf{22.} $0.05 \times 10^7/3600 = \qty{139}{m^3/s}$; $v = 139/40 = \qty{3.5}{m/s}$.

\textbf{23.} $D_V = \qty{83}{cm^3/s}$: aorta \qty{0.33}{m/s}; kapiler
\qty{0.33}{mm/s}; \qty{3}{s} tiap milimeter.

\textbf{24.} $S\,\dd h/\dd t = -s\sqrt{2gh}$: $\sqrt h = \sqrt{h_0} -
\frac sS\sqrt{g/2}\,t$, kosong pada $T = (S/s)\sqrt{2h_0/g} = 10^4 \times 0.45 =
\qty{4500}{s} \approx \qty{75}{min}$.

\textbf{25.} Bak: $r \sim \qty{0.5}{m}$, $v \sim \qty{1}{cm/s}$, $\Gamma =
\qty{3e-2}{m^2/s}$; puting beliung: \qty{60}{m}, \qty{80}{m/s}, $\omega = \qty{2.7}{s^{-1}}$,
$\Gamma = \qty{3e4}{m^2/s}$; siklon: \qty{40}{km}, \qty{50}{m/s}, $\omega =
\qty{2.5e-3}{s^{-1}}$, $\Gamma = \qty{1.3e7}{m^2/s}$. Pada masing-masingnya, momentum
sudut tiap satuan massa $rv_\theta$ fluida yang memusat itu kekal:
aliran masuknyalah yang memutarnya makin cepat.
\end{solution}
